\documentclass[10pt,a4paper]{amsart}
\usepackage{amsmath,amssymb,amsthm,graphicx}
% Retained minimal preamble for the Stacks Project chapter sites-cohomology.tex (Cohomology on Sites).
% Verbatim excerpt of the chapter preamble: Xy-pic setup (native preamble lines 14-18) and the native
% theorem declarations with the mathematical macros (native preamble lines 154-217). Copied mechanically;
% nothing added, removed or redefined. The wrapper class and the amsmath/amssymb/amsthm/graphicx base
% packages are supplied by the assembler, so the class/externaldocument part of the native preamble is omitted.
% For commutative diagrams we use Xy-pic
\usepackage[all]{xy}

% We use 2cell for 2-commutative diagrams.
\xyoption{2cell}
% Theorem environments.
%
\theoremstyle{plain}
\newtheorem{theorem}[subsection]{Theorem}
\newtheorem{proposition}[subsection]{Proposition}
\newtheorem{lemma}[subsection]{Lemma}

\theoremstyle{definition}
\newtheorem{definition}[subsection]{Definition}
\newtheorem{example}[subsection]{Example}
\newtheorem{exercise}[subsection]{Exercise}
\newtheorem{situation}[subsection]{Situation}

\theoremstyle{remark}
\newtheorem{remark}[subsection]{Remark}
\newtheorem{remarks}[subsection]{Remarks}

\numberwithin{equation}{subsection}

% Macros
%
\def\lim{\mathop{\mathrm{lim}}\nolimits}
\def\colim{\mathop{\mathrm{colim}}\nolimits}
\def\Spec{\mathop{\mathrm{Spec}}}
\def\Hom{\mathop{\mathrm{Hom}}\nolimits}
\def\Ext{\mathop{\mathrm{Ext}}\nolimits}
\def\SheafHom{\mathop{\mathcal{H}\!\mathit{om}}\nolimits}
\def\SheafExt{\mathop{\mathcal{E}\!\mathit{xt}}\nolimits}
\def\Sch{\mathit{Sch}}
\def\Mor{\mathop{\mathrm{Mor}}\nolimits}
\def\Ob{\mathop{\mathrm{Ob}}\nolimits}
\def\Sh{\mathop{\mathit{Sh}}\nolimits}
\def\NL{\mathop{N\!L}\nolimits}
\def\CH{\mathop{\mathrm{CH}}\nolimits}
\def\proetale{{pro\text{-}\acute{e}tale}}
\def\etale{{\acute{e}tale}}
\def\QCoh{\mathit{QCoh}}
\def\Ker{\mathop{\mathrm{Ker}}}
\def\Im{\mathop{\mathrm{Im}}}
\def\Coker{\mathop{\mathrm{Coker}}}
\def\Coim{\mathop{\mathrm{Coim}}}

% Boxtimes
%
\DeclareMathSymbol{\boxtimes}{\mathbin}{AMSa}{"02}

%
% Macros for moduli stacks/spaces
%
\def\QCohstack{\mathcal{QC}\!\mathit{oh}}
\def\Cohstack{\mathcal{C}\!\mathit{oh}}
\def\Spacesstack{\mathcal{S}\!\mathit{paces}}
\def\Quotfunctor{\mathrm{Quot}}
\def\Hilbfunctor{\mathrm{Hilb}}
\def\Curvesstack{\mathcal{C}\!\mathit{urves}}
\def\Polarizedstack{\mathcal{P}\!\mathit{olarized}}
\def\Complexesstack{\mathcal{C}\!\mathit{omplexes}}
% \Pic is the operator that assigns to X its picard group, usage \Pic(X)
% \Picardstack_{X/B} denotes the Picard stack of X over B
% \Picardfunctor_{X/B} denotes the Picard functor of X over B
\def\Pic{\mathop{\mathrm{Pic}}\nolimits}
\def\Picardstack{\mathcal{P}\!\mathit{ic}}
\def\Picardfunctor{\mathrm{Pic}}
\def\Deformationcategory{\mathcal{D}\!\mathit{ef}}
\usepackage{cleveref}
\crefname{originallemma}{original lemma}{original lemmas}
\Crefname{originallemma}{Original lemma}{Original lemmas}
\makeatletter
\newlabel{derived-proposition-homotopy-category-triangulated}{{014S}{source}{}{sourcepointer-1}{}}
\newlabel{derived-proposition-homotopy-category-triangulated@cref}{{[originallemma][0][]014S}{[1][0][]source}{}{}{}}
\makeatother
\crefname{originallemma}{original lemma}{original lemmas}
\Crefname{originallemma}{Original lemma}{Original lemmas}
\makeatletter
\newlabel{lemma-K-flat-resolution}{{06YS}{source}{}{sourcepointer-2}{}}
\newlabel{lemma-K-flat-resolution@cref}{{[originallemma][0][]06YS}{[1][0][]source}{}{}{}}
\makeatother
\crefname{originallemma}{original lemma}{original lemmas}
\Crefname{originallemma}{Original lemma}{Original lemmas}
\makeatletter
\newlabel{lemma-K-flat-two-out-of-three}{{07A3}{source}{}{sourcepointer-3}{}}
\newlabel{lemma-K-flat-two-out-of-three@cref}{{[originallemma][0][]07A3}{[1][0][]source}{}{}{}}
\makeatother
\crefname{originallemma}{original lemma}{original lemmas}
\Crefname{originallemma}{Original lemma}{Original lemmas}
\makeatletter
\newlabel{derived-lemma-third-isomorphism-triangle}{{014A}{source}{}{sourcepointer-4}{}}
\newlabel{derived-lemma-third-isomorphism-triangle@cref}{{[originallemma][0][]014A}{[1][0][]source}{}{}{}}
\makeatother
\crefname{originallemma}{original lemma}{original lemmas}
\Crefname{originallemma}{Original lemma}{Original lemmas}
\makeatletter
\newlabel{derived-lemma-improve-distinguished-triangle-homotopy}{{0G6C}{source}{}{sourcepointer-5}{}}
\newlabel{derived-lemma-improve-distinguished-triangle-homotopy@cref}{{[originallemma][0][]0G6C}{[1][0][]source}{}{}{}}
\makeatother
\title{Factoring a map out of a K-flat complex through a K-flat quasi-isomorphism}
\author{The Stacks Project Authors}
\date{stacks-project:sites-cohomology.tex:local-snapshot-20261011; GNU FDL 1.2 or later}
\begin{document}
\maketitle
\noindent\textit{Editorial note.} All delivered mathematical text is the author's own, copied from the pinned Stacks Project chapter file sites-cohomology.tex; no statement or proof step was rewritten, reordered or removed.\par\medskip
\noindent\textit{Editorial note.} The proof is the author's complete argument for all clauses of the statement, including the preservation of flatness of the terms, and nothing is delegated.\par\medskip
\noindent\textit{Editorial note.} Every labelled reference occurring inside the counted spans is carried in assembly.external\_references as a non-load-bearing navigation pointer (same-chapter ones at their definition site, cross-chapter ones at the invocation inside the counted span), so no load-bearing mathematics is replaced by a stub. No bibliography citation occurs inside the counted spans.\par\medskip
\noindent\textit{Editorial note.} The checked chapter list for same-name collisions is the set of stacks chapters staged in this repository: chow, cohomology, divisors, etale-cohomology, flat, more-algebra, more-morphisms, morphisms, spaces-more-morphisms, varieties together with this card's chapter sites-cohomology.\par\medskip
\noindent\textit{Editorial note.} The preamble is the chapter's own retained minimal preamble (sources/stacks-sites-cohomology/retained-preamble.tex), a verbatim excerpt of the native preamble: the Xy-pic setup and the native theorem declarations with the mathematical macros.\par\medskip
\noindent\textit{Allowed standard prerequisites (editorial).}\begin{itemize}
\item K-flat complex of O-modules and the flatness of its terms. Same chapter, the definition of K-flatness and the criterion by acyclicity of total tensor products. Supplies the hypothesis on K\textasciicircum{}bullet and the property required of the auxiliary complex N\textasciicircum{}bullet.
\item Triangulated structure of the homotopy category K(Mod(O)). Same chapter, distinguished triangles in the homotopy category and the long exact sequence of homotopy classes they induce. Supplies the argument by which the map a is replaced up to homotopy by a map into a quasi-isomorphic complex.
\item Existence of K-flat resolutions with flat terms. Same chapter, the construction of a K-flat complex mapping quasi-isomorphically to a given complex, respecting flatness of the terms. Supplies the construction of N\textasciicircum{}bullet, the map c and their properties in the second part of the statement.
\end{itemize}\medskip
\ifdefined\hypertarget\hypertarget{sourcepointer-1}{}\fi
\noindent\textit{Source reference.} Original lemma 014S: sources/stacks-sites-cohomology/pinned-original-source.tex, source lines 3438--3438. Navigation pointer invoked at proof line 3438; the counted span contains the pointer only and the cited result is not reproduced inside it, so the pointer is registered as a non-load-bearing navigation stub in the established form for Stacks material. No bibliography citation is involved.\par\medskip
\ifdefined\hypertarget\hypertarget{sourcepointer-2}{}\fi
\noindent\textit{Source reference.} Original lemma 06YS: sources/stacks-sites-cohomology/pinned-original-source.tex, source lines 3207--3207. Navigation pointer invoked at proof line 3444; the counted span contains the pointer only and the cited result is not reproduced inside it, so the pointer is registered as a non-load-bearing navigation stub in the established form for Stacks material. No bibliography citation is involved.\par\medskip
\ifdefined\hypertarget\hypertarget{sourcepointer-3}{}\fi
\noindent\textit{Source reference.} Original lemma 07A3: sources/stacks-sites-cohomology/pinned-original-source.tex, source lines 3081--3081. Navigation pointer invoked at proof line 3451; the counted span contains the pointer only and the cited result is not reproduced inside it, so the pointer is registered as a non-load-bearing navigation stub in the established form for Stacks material. No bibliography citation is involved.\par\medskip
\ifdefined\hypertarget\hypertarget{sourcepointer-4}{}\fi
\noindent\textit{Source reference.} Original lemma 014A: sources/stacks-sites-cohomology/pinned-original-source.tex, source lines 3473--3473. Navigation pointer invoked at proof line 3473; the counted span contains the pointer only and the cited result is not reproduced inside it, so the pointer is registered as a non-load-bearing navigation stub in the established form for Stacks material. No bibliography citation is involved.\par\medskip
\ifdefined\hypertarget\hypertarget{sourcepointer-5}{}\fi
\noindent\textit{Source reference.} Original lemma 0G6C: sources/stacks-sites-cohomology/pinned-original-source.tex, source lines 3478--3478. Navigation pointer invoked at proof line 3478; the counted span contains the pointer only and the cited result is not reproduced inside it, so the pointer is registered as a non-load-bearing navigation stub in the established form for Stacks material. No bibliography citation is involved.\par\medskip
\begin{lemma}
\label{lemma-factor-through-K-flat}
Let $(\mathcal{C}, \mathcal{O})$ be a ringed site.
Let $a : \mathcal{K}^\bullet \to \mathcal{L}^\bullet$ be a map of complexes
of $\mathcal{O}$-modules. If $\mathcal{K}^\bullet$ is K-flat, then
there exist a complex $\mathcal{N}^\bullet$ and maps of complexes
$b : \mathcal{K}^\bullet \to \mathcal{N}^\bullet$
and $c : \mathcal{N}^\bullet \to \mathcal{L}^\bullet$ such that
\begin{enumerate}
\item $\mathcal{N}^\bullet$ is K-flat,
\item $c$ is a quasi-isomorphism,
\item $a$ is homotopic to $c \circ b$.
\end{enumerate}
If the terms of $\mathcal{K}^\bullet$ are flat, then we may choose
$\mathcal{N}^\bullet$, $b$, and $c$
such that the same is true for $\mathcal{N}^\bullet$.
\end{lemma}

\begin{proof}
We will use that the homotopy category $K(\textit{Mod}(\mathcal{O}))$
is a triangulated category, see Derived Categories, Proposition
\ref{derived-proposition-homotopy-category-triangulated}.
Choose a distinguished triangle
$\mathcal{K}^\bullet \to \mathcal{L}^\bullet \to
\mathcal{C}^\bullet \to \mathcal{K}^\bullet[1]$.
Choose a quasi-isomorphism $\mathcal{M}^\bullet \to \mathcal{C}^\bullet$ with
$\mathcal{M}^\bullet$ K-flat with flat terms, see
Lemma \ref{lemma-K-flat-resolution}.
By the axioms of triangulated categories,
we may fit the composition
$\mathcal{M}^\bullet \to \mathcal{C}^\bullet \to \mathcal{K}^\bullet[1]$
into a distinguished triangle
$\mathcal{K}^\bullet \to \mathcal{N}^\bullet \to
\mathcal{M}^\bullet \to \mathcal{K}^\bullet[1]$.
By Lemma \ref{lemma-K-flat-two-out-of-three} we see that
$\mathcal{N}^\bullet$ is K-flat.
Again using the axioms of triangulated categories,
we can choose a map $\mathcal{N}^\bullet \to \mathcal{L}^\bullet$ fitting into
the following morphism of distinguished triangles
$$
\xymatrix{
\mathcal{K}^\bullet \ar[r] \ar[d] &
\mathcal{N}^\bullet \ar[r] \ar[d] &
\mathcal{M}^\bullet \ar[r] \ar[d] &
\mathcal{K}^\bullet[1] \ar[d] \\
\mathcal{K}^\bullet \ar[r] &
\mathcal{L}^\bullet \ar[r] &
\mathcal{C}^\bullet \ar[r] &
\mathcal{K}^\bullet[1]
}
$$
Since two out of three of the arrows are quasi-isomorphisms, so is
the third arrow $\mathcal{N}^\bullet \to \mathcal{L}^\bullet$
by the long exact sequences
of cohomology associated to these distinguished triangles
(or you can look at the image of this diagram in $D(\mathcal{O})$ and use
Derived Categories, Lemma \ref{derived-lemma-third-isomorphism-triangle}
if you like). This finishes the proof of (1), (2), and (3).
To prove the final assertion, we may choose $\mathcal{N}^\bullet$ such that
$\mathcal{N}^n \cong \mathcal{M}^n \oplus \mathcal{K}^n$, see
Derived Categories, Lemma
\ref{derived-lemma-improve-distinguished-triangle-homotopy}.
Hence we get the desired flatness
if the terms of $\mathcal{K}^\bullet$ are flat.
\end{proof}

\end{document}
